\part{Fundamentos}

\section{Mecânica Básica}

Vanguarda 44 utiliza rolagens de dados d6 para quase todas as ações do jogo. O mecanismo central está centrado nos Pontos de Comando (PC).

\subsection{Pontos de Comando}

Cada NCO (sargento) ou oficial de campo (tenente, capitão ou major) no campo de batalha representa uma fonte de Pontos de Comando, que variam de acordo com o posto militar:

\begin{center}
\begin{tabular}{lcc}
\toprule
Posto & Suporte de Base & PC Gerados \\
\midrule
Sargento / Líder de Esquadrão & Amarelo Ocre (Mostarda) & 1 \\
Tenente & Vermelho Escuro (Vinho) & 2 \\
Capitão & Laranja & 3 \\
Major & --- & 4 \\
General & --- & 6 \\
\bottomrule
\end{tabular}
\end{center}

\begin{nota}
Nem todo esquadrão possui um sargento ou oficial de campo com PC. Times sem líder próprio (como uma equipe de Metralhadora Pesada) precisam que um oficial gaste 1 PC extra para ativá-los --- veja a Parte VI, Estrutura de Comando.
\end{nota}

A identificação visual por cores dos suportes de base é universal entre os exércitos: mesmo que cada nação tenha uniformes diferentes, um suporte vinho é sempre um Tenente (2 PC) e um suporte mostarda é sempre um Sargento (1 PC). Isso permite que ambos os jogadores leiam a mesa rapidamente sem precisar consultar fichas.

\subsection{Saco de Dados de Ordem e Ativação}

\textbf{Saco de Dados de Ordem}: um saco ou recipiente contém um Dado de Ordem para cada unidade no campo de batalha.

\textbf{Ativação}: os jogadores se revezam sorteando um único Dado de Ordem do saco. O jogador cuja cor é sorteada ativa uma de suas unidades que ainda não recebeu ordem.

Cada Ponto de Comando disponível pode ser usado para ativar um esquadrão em campo.

\subsection{Ordens}

Quando ativada, uma unidade tenta executar uma das seguintes ordens:

\begin{itemize}[leftmargin=1.2cm]
  \item \textbf{Avançar}: mover e depois atirar com eficácia reduzida (-1 para acertar).
  \item \textbf{Correr}: mover o dobro da distância, mas não pode atirar.
  \item \textbf{Atirar}: ficar parado ou pivotar e atirar com eficácia total.
  \item \textbf{Reagrupar (Rally)}: tentar remover Marcadores de Stress (veja Parte III, Moral e Supressão).
  \item \textbf{Emboscada (Ambush)}: esperar que uma unidade inimiga se mova ou atire para então reagir.
  \item \textbf{Abaixar (Down)}: ficar prostrado para defesa máxima.
\end{itemize}

\subsection{Teste de Ordem}

Para executar uma ordem com sucesso, a unidade deve passar em um Teste de Ordem: role 1D6 contra um Número Alvo (TN).

\textbf{TN base}: 3+ para uma unidade padrão sem Marcadores de Stress.

\textbf{Modificador}: +1 ao TN para cada Marcador de Stress que a unidade possua.

\textbf{Resultado}: se a rolagem atingir ou exceder o TN, a ordem é executada. Se falhar, o dado da unidade é imediatamente colocado na posição Abaixar, e a unidade não pode ser ativada novamente neste turno.

\section{Identificação de Comando (Suportes de Base)}

Para garantir a fluidez do jogo e evitar confusões no caos do campo de batalha, Vanguarda 44 utiliza um sistema de identificação visual por cores nos suportes das bases. Embora cada exército possua seus próprios uniformes e insígnias, as cores dos suportes de comando são universais, permitindo que ambos os jogadores identifiquem instantaneamente as fontes de Pontos de Comando na mesa.

\textbf{Vermelho Escuro (Vinho)}: identifica o Tenente (Oficial) --- gera 2 PC por turno.

\textbf{Amarelo Ocre (Mostarda)}: identifica os Sargentos (NCOs) de esquadrão --- gera 1 PC por turno.

Essa padronização visual é fundamental para a leitura da mesa, permitindo que os jogadores foquem nas decisões estratégicas sem interromper o adversário para identificar quem é o líder de cada unidade.
